% GENERATED FILE. Edit canonical lessons, then run npm run generate:books.
% canonical-source-hash: fnv1a64:65653e2907d844f1
% canonical-lessons: ZH-C83-shijian, ZH-C83-zao, ZH-C83-mashang, ZH-C83-laibuji, ZH-C83-chidao

\chapter{Running Late}
\label{ch:zh-running-late}
\hlchaptermodality{83}

\begin{chapteropening}
\textbf{By the end of this chapter:} \emph{I can ask if somebody is free, and say early, right away, too late and late.}

Complete the last lesson of chapter 83: I can ask if somebody is free, and say early, right away, too late and late.
\end{chapteropening}

\section[shíjiān]{shíjiān — time; free time}

\label{lesson:ZH-C83-shijian}



\begin{quote}
Before the new one: say the Mandarin for the middle, then the Mandarin for opposite.
\end{quote}

\subsection*{You'll want to know: shíjiān}
\emph{shíjiān} — \textquotedblleft{}time; free time\textquotedblright{}.

\begin{quote}
\emph{nǐ míngtiān yǒu shíjiān ma} — Are you free tomorrow?
\end{quote}

Its characters come later; for now it is a word for the ear.

\begin{cousinweb}
\emph{shíjiān} is built from shí 'time' and jiān 'a gap', both inherited from Middle Chinese, \emph{dzyi} and \emph{kean} (Guangyun, in Baxter's transcription); the dictionary glosses the first with another word for a time, and lists the second under an older form of the character.
\end{cousinweb}

\begin{grammarlens}[title={What you've built}]
One more word for everyday answers.
\end{grammarlens}

\subsection*{Your turn}
\begin{itemize}
\raggedright
  \item \emph{Say it:} \emph{shíjiān}
  \item \emph{Say it:} \emph{shíjiān}, once more
  \item \emph{Say it:} \emph{nǐ míngtiān yǒu shíjiān ma} — Are you free tomorrow?
  \item \emph{Recall:} say the Mandarin for the middle, then the Mandarin for opposite, then say \emph{shíjiān} again
\end{itemize}

\begin{tcolorbox}[breakable,colback=teal!4,colframe=teal!35!black,title={Before you move on}]
What does shíjiān mean? (\textquotedblleft{}Time; free time\textquotedblright{}.) Say it once more.
\end{tcolorbox}
\section[\texorpdfstring{\zh{早} (zǎo)}{早 (zǎo)}]{\zh{早} (zǎo) — early}

\label{lesson:ZH-C83-zao}



\begin{quote}
Before the new one: say the Mandarin for opposite, then the Mandarin for time to spare.
\end{quote}

\subsection*{You'll want to know: \zh{早}}
\textbf{\zh{早}} — \emph{zǎo} — \textquotedblleft{}early\textquotedblright{}.

\begin{quote}
\textbf{\zh{明天的火车很早}} \emph{míngtiān de huǒchē hěn zǎo} — Tomorrow's train leaves early.
\end{quote}

\begin{cousinweb}
\textbf{\zh{早}} (\emph{zǎo}) is inherited from Middle Chinese \emph{tsawX} (Guangyun, in Baxter's transcription), glossed 'morning'. Early is not given there; from morning to early is the likely step, but it is not settled.
\end{cousinweb}

\begin{grammarlens}[title={What you've built}]
One more word for everyday answers.
\end{grammarlens}

\subsection*{Your turn}
\begin{itemize}
\raggedright
  \item \emph{Say it:} \emph{zǎo}
  \item \emph{Say it:} \emph{zǎo}, once more
  \item \emph{Say it:} \emph{míngtiān de huǒchē hěn zǎo} — Tomorrow's train leaves early.
  \item \emph{Recall:} say the Mandarin for opposite, then the Mandarin for time to spare, then say \emph{zǎo} again
\end{itemize}

\begin{tcolorbox}[breakable,colback=teal!4,colframe=teal!35!black,title={Before you move on}]
What does \zh{早} mean? (\textquotedblleft{}Early\textquotedblright{}.) Say it once more.
\end{tcolorbox}
\section[\texorpdfstring{\zh{马上} (mǎshàng)}{马上 (mǎshàng)}]{\zh{马上} (mǎshàng) — at once; right away}

\label{lesson:ZH-C83-mashang}



\begin{quote}
Before the new one: say the Mandarin for time to spare, then the Mandarin for early.
\end{quote}

\subsection*{You'll want to know: \zh{马上}}
\textbf{\zh{马上}} — \emph{mǎshàng} — \textquotedblleft{}at once; right away\textquotedblright{}.

\begin{quote}
\textbf{\zh{家里人马上来}} \emph{jiā lǐ rén mǎshàng lái} — The family will be here any minute.
\end{quote}

\begin{cousinweb}
\textbf{\zh{马上}} (\emph{mǎshàng}) is built from \zh{马} 'horse' and \zh{上} 'on', both inherited from Middle Chinese, \emph{maeX} and \emph{dzyangX} (Guangyun, in Baxter's transcription). 'On horseback', so already on the way, is the usual explanation, but it is not settled.
\end{cousinweb}

\begin{grammarlens}[title={What you've built}]
One more word for everyday answers.
\end{grammarlens}

\subsection*{Your turn}
\begin{itemize}
\raggedright
  \item \emph{Say it:} \emph{mǎshàng}
  \item \emph{Say it:} \emph{mǎshàng}, once more
  \item \emph{Say it:} \emph{jiā lǐ rén mǎshàng lái} — The family will be here any minute.
  \item \emph{Recall:} say the Mandarin for time to spare, then the Mandarin for early, then say \emph{mǎshàng} again
\end{itemize}

\begin{tcolorbox}[breakable,colback=teal!4,colframe=teal!35!black,title={Before you move on}]
What does \zh{马上} mean? (\textquotedblleft{}At once; right away\textquotedblright{}.) Say it once more.
\end{tcolorbox}
\section[láibují]{láibují — to run out of time}

\label{lesson:ZH-C83-laibuji}



\begin{quote}
Before the new one: say the Mandarin for early, then the Mandarin for right away.
\end{quote}

\subsection*{You'll want to know: láibují}
\emph{láibují} — \textquotedblleft{}to run out of time\textquotedblright{}.

\begin{quote}
\emph{xiànzài qù láibují le} — It's too late to go now.
\end{quote}

Its characters come later; for now it is a word for the ear.

Its middle syllable loses its tone: láibují, with bu light between two full syllables.

\begin{cousinweb}
\emph{láibují} is built from \zh{来} 'to come', \zh{不} 'not' and jí 'to reach', all inherited from Middle Chinese, \emph{loj}, \emph{pjut} and \emph{gip} (Guangyun, in Baxter's transcription): coming, but not reaching in time.
\end{cousinweb}

\begin{grammarlens}[title={What you've built}]
One more word for everyday answers.
\end{grammarlens}

\subsection*{Your turn}
\begin{itemize}
\raggedright
  \item \emph{Say it:} \emph{láibují}
  \item \emph{Say it:} \emph{láibují}, once more
  \item \emph{Say it:} \emph{xiànzài qù láibují le} — It's too late to go now.
  \item \emph{Recall:} say the Mandarin for early, then the Mandarin for right away, then say \emph{láibují} again
\end{itemize}

\begin{tcolorbox}[breakable,colback=teal!4,colframe=teal!35!black,title={Before you move on}]
What does láibují mean? (\textquotedblleft{}To run out of time\textquotedblright{}.) Say it once more.
\end{tcolorbox}
\section[chídào]{chídào — to be late; to arrive late}

\label{lesson:ZH-C83-chidao}



\begin{quote}
Before the new one: say the Mandarin for right away, then the Mandarin for to run out of time.
\end{quote}

\subsection*{You'll want to know: chídào}
\emph{chídào} — \textquotedblleft{}to be late; to arrive late\textquotedblright{}.

\begin{quote}
\emph{duìbuqǐ wǒ chídào le} — Sorry I'm late.
\end{quote}

Its characters come later; for now it is a word for the ear.

\begin{cousinweb}
\emph{chídào} is built from chí 'slow' and dào 'to arrive', both inherited from Middle Chinese, \emph{drij} and \emph{tawH} (Guangyun, in Baxter's transcription).
\end{cousinweb}

\begin{grammarlens}[title={What you've built}]
One more word for everyday answers.
\end{grammarlens}

\subsection*{Your turn}
\begin{itemize}
\raggedright
  \item \emph{Say it:} \emph{chídào}
  \item \emph{Say it:} \emph{chídào}, once more
  \item \emph{Say it:} \emph{duìbuqǐ wǒ chídào le} — Sorry I'm late.
  \item \emph{Recall:} say the Mandarin for right away, then the Mandarin for to run out of time, then say \emph{chídào} again
\end{itemize}

\begin{tcolorbox}[breakable,colback=teal!4,colframe=teal!35!black,title={Before you move on}]
What does chídào mean? (\textquotedblleft{}To be late; to arrive late\textquotedblright{}.) Say it once more.
\end{tcolorbox}
